\chapter{Discusión y limitaciones}
\label{ch:discusion}
% Fuente: FINAL_HOLDOUT_V2_RESULTS.md §§8--10 y Discussion-ready statements.
\section{Interpretación de los resultados primarios}
\pendiente{Recordar brevemente qué mide F1 antes de interpretar cifras.
Explicar su significado respecto a los objetivos antes de introducir
el diagnóstico. Separar detección de entidades y atributos sobre pares
emparejados; no atribuir toda la salida final al modelo Gemini.}
\section{Diagnóstico posterior de las discrepancias}
\pendiente{POST-HOC DIAGNOSTIC — NOT PRIMARY SCORING. Utilizar únicamente
el diagnóstico aprobado: alias, división de eventos, dependencia jerárquica,
evidencia de títulos y canonicalización. Distinguir error del modelo, efecto
determinista, limitación del matching y discrepancia de representación;
no confundirlos con errores de agrupación longitudinal de Gold.
No crear nuevos matches.}
\pendiente{FIGURA F08 — Cascada de matching
\par Etiqueta prevista: \path{fig:cascada-matching}. Etiquetar POST-HOC DIAGNOSTIC — NOT PRIMARY
SCORING. La ausencia de emparejamiento de un activo impide satisfacer el
conjunto completo exigido para su evento; esto puede dejar actuaciones y
localizaciones sin match. El emparejamiento del padre es necesario, no
suficiente. Si ayudan, usar los conteos diagnósticos ya fijados: activos
37 truth / 39 predichos / 20 matches; eventos 30 / 38 / 13; actuaciones
36 actuales truth / 58 predichas / 1 match actual. No convertirlos en una
métrica nueva ni presentar el diagrama como cambio del scoring. Referenciar
desde esta sección y mantenerlo separado de los resultados primarios.}
\section{Validez y límites de la evaluación}
\pendiente{Describir tamaño y selección del holdout, una anotadora principal,
QA realmente documentado, denominadores pequeños, P0 sin positivos
emparejados, evidencia y error terminal conservado. No inferir equivalencia
semántica por comparación aproximada ni generalizar atributos 1/1.}
\section{Límites del corpus y utilidad del sistema}
\pendiente{Relacionar alcance institucional y temporal, recuperación
conservadora, incertidumbre territorial y dependencia del proveedor con lo
que puede consultarse y concluirse. Separar utilidad técnica demostrada de
usabilidad o calidad longitudinal no evaluadas mediante este holdout.}
\pendiente{TABLA — Limitaciones: alcance, efecto sobre la interpretación
y evidencia disponible. Separar corpus, extracción, grouping y evaluación;
no repetir como lista de cifras las tablas de resultados.}
